\begin{statementbox}
\textbf{条件：} 直线 $r$ 关于镜面直线 $l$ 反射为 $r'$。原斜率记作 $k$，像斜率记作 $k'$；谈有限斜率时，相应直线必须不竖直。

\begin{center}
% Equal units on both axes. Blue original, orange image, teal mirror.
\begin{tikzpicture}[x=1cm,y=1cm]
\draw[p055 frame] (0,0) rectangle (13.8,6.2);
\node[anchor=west,font=\small\bfseries] at (.4,5.7) {同一条原线：镜面的方向决定怎样变斜率};
\begin{scope}[shift={(3.5,2.9)},scale=.95]
  \draw[p055 axis] (-2.8,0)--(2.8,0) node[right] {$x$};
  \draw[p055 mirror] (0,-1.2)--(0,2.2);
  \node[teal!70!black,anchor=east] at (-.12,2) {$x=0$};
  \draw[p055 original] (-2.6,-1.3)--(2.6,1.3);
  \draw[p055 image] (-2.6,1.3)--(2.6,-1.3);
  \fill[CStatement] (2,1) circle (2pt) node[above=12pt] {$(2,1)$};
  \fill[orange!85!black] (-2,1) circle (2pt) node[above=12pt] {$(-2,1)$};
  \draw[p055 helper] (-2,1)--(2,1);
\end{scope}
\begin{scope}[shift={(10.25,2.35)},scale=.95]
  \draw[p055 axis] (-.5,0)--(2.9,0) node[right] {$x$};
  \draw[p055 mirror] (-.4,-.4)--(2.5,2.5) node[above] {$y=x$};
  \draw[p055 original] (-.5,-.25)--(2.8,1.4);
  \draw[p055 image] (-.25,-.5)--(1.4,2.8);
  \fill[CStatement] (2,1) circle (2pt) node[below right=3pt] {$(2,1)$};
  \fill[orange!85!black] (1,2) circle (2pt) node[left] {$(1,2)$};
  \draw[p055 helper] (2,1)--(1,2);
\end{scope}
\node[align=center] at (3.5,.85) {左右翻转：$\frac12\mapsto-\frac12$\\水平镜面同样只翻转一个增量};
\node[align=center] at (10.3,.85) {横纵交换：$\frac12\mapsto2$\\另一种 $45^\circ$ 镜面同时反号，商不变};
\end{tikzpicture}

\end{center}

\textbf{两类常用规则（先辨镜面方向）：}
\[
\boxed{\begin{array}{c|c|c}
\text{镜面}&(\Delta x,\Delta y)\text{ 变成}&\text{像斜率}\\ \hline
x=a&(-\Delta x,\Delta y)&-k\\
y=b&(\Delta x,-\Delta y)&-k\\
y=x+c&(\Delta y,\Delta x)&1/k\\
y=-x+c&(-\Delta y,-\Delta x)&1/k
\end{array}}
\]
其中 $\Delta x,\Delta y$ 是原线上两个不同点的坐标差，$k=\Delta y/\Delta x$。倒数规则限于 $k\ne0$ 且有限；水平线经 $45^\circ$ 镜面变竖直，竖直线则变水平。

\medskip
\textbf{任意非竖直镜面：} 若 $l:y=mx+c$，$k$ 有限，则
\[
\boxed{k'=\frac{2m+(m^2-1)k}{1-m^2+2mk}}\qquad(1-m^2+2mk\ne0).
\]
分母为 $0$ 表示像直线竖直；镜面 $x=a$ 则直接使用 $-k$。镜面位置参数 $a,b,c$ 不影响上述方向变换，但影响像直线的位置。

\medskip
\textbf{方向之外还要核对：} 两条不同的相交直线若交点 $P$ 在镜面上，再核对方向即可判断是否互为反射。用于三角形时，还须确认镜面平分的是两条边\emph{射线}的内角；仅有斜率关系不能排除外角平分线。
\end{statementbox}
